%% ---------------------------------------------------------------- title page
\thispagestyle{empty}
\begin{center}
\vspace*{1.2in}
{\sans\bfseries\fontsize{40}{46}\selectfont <<TITLE>>\par}
\vspace{0.25in}
{\sans\fontsize{18}{24}\selectfont <<SUBTITLE>>\par}
\vspace{0.12in}
{\sans\large\color{mid}<<TAGLINE>>\par}
\vspace{0.5in}
{\color{ink}\rule{2.2in}{1.2pt}}\par
\vspace{0.5in}
\begin{minipage}{4.6in}\centering\large
Key concepts and worked examples for every topic\\[2pt]
A diagnostic test with a personal study plan\\[2pt]
<<NPRACTICE>> drill problems in <<NCHAPTERS>> chapters\\[2pt]
4 full-length math practice tests\\[2pt]
Every answer checked by computer algebra
\end{minipage}
\vfill
{\sans\Large <<AUTHOR>>\par}
\vspace{0.4in}
\end{center}
\clearpage

%% ---------------------------------------------------------------- copyright
\thispagestyle{empty}
\vspace*{\fill}
{\small
\textbf{<<TITLE>>: <<SUBTITLE>>}\par
<<EDITION>>, <<YEAR>>\par
\medskip
Copyright \textcopyright{} <<YEAR>> <<AUTHOR>>. All rights reserved.\par
No part of this book may be reproduced, stored in a retrieval system, or
transmitted in any form or by any means without the prior written permission
of the author, except for brief quotations in reviews.\par
\medskip
ASVAB is a registered trademark of the United States Department of Defense,
which was not involved in the production of, and does not endorse, this book.
This book is an independent study aid. It is not affiliated with the
Department of Defense, any branch of the U.S.\ Armed Forces, or any
recruiting organization.\par
\medskip
Test formats, question counts, time limits, and qualifying scores are set by
the Department of Defense and the individual services and can change. Always
confirm current details with your recruiter or at officialasvab.com. The
author makes no guarantee of any particular test score.\par
\medskip
The practice problems in this book were generated and their answers verified
with the SymPy computer algebra system; every answer was computed exactly and
then confirmed by a second, independent calculation.\par
}
\clearpage

\tableofcontents
\clearpage

%% ---------------------------------------------------------------- how to use
\chapter*{How to Use This Book}
\addcontentsline{toc}{chapter}{How to Use This Book}
\markboth{How to Use This Book}{How to Use This Book}

The math on the ASVAB is mostly middle-school and early high-school
math---percents, fractions, equations, areas, word problems---asked quickly
and without a calculator. What raises scores is \emph{practice}: seeing each
kind of question often enough that you recognize it instantly and know the
first step.

This workbook gives you that practice: \textbf{1,000 problems}, each with a
complete step-by-step solution.

\section*{The plan}
\begin{enumerate}[leftmargin=1.6em,itemsep=0.4em]
\item \textbf{Take the Diagnostic Test} (page~\pageref{diag}). Thirty
  questions, about 35 minutes, no calculator. The answer key tells you which
  chapter each question comes from.
\item \textbf{Study the chapters you need.} Start with the topics you missed
  on the diagnostic. Each chapter has three parts:
  \begin{itemize}
  \item \emph{Key Concepts}---the rules, formulas, a worked example, a
    no-calculator tip, and the traps test writers love.
  \item \emph{Practice Set}---warm-up problems first, then test-level
    problems, then challenge problems.
  \item \emph{Answers \& Explanations}---a quick answer key plus a full
    solution for every problem.
  \end{itemize}
\item \textbf{Take the four Practice Tests} under real time limits once you
  have worked through the chapters. Each has an Arithmetic Reasoning section
  (30 questions, 36 minutes) and a Mathematics Knowledge section
  (25 questions, 24 minutes), the format of the paper ASVAB.
\item \textbf{Review every miss.} The answer keys of the tests list the
  chapter for every question, so you always know where to go back.
\end{enumerate}

\section*{Reading the solutions}
Every solution is broken into numbered steps that you can follow with pencil
and paper. Many solutions also include:
\begin{itemize}
\item \textbf{TIP}---a shortcut, an estimate, or a way to check your answer.
\item \textbf{WHY NOT}---what went wrong if you picked one of the tempting
  wrong answers. The wrong choices in this book are not random numbers; each
  comes from a real, common mistake (dividing by the wrong number, forgetting
  to convert minutes to hours, adding percents that should be multiplied).
  If your answer matches one of them, the explanation tells you exactly which
  mistake you made.
\end{itemize}

\section*{Good habits}
\begin{itemize}
\item \textbf{No calculator.} The ASVAB does not allow one, so don't practice
  with one. Every problem here is designed to be done by hand.
\item \textbf{Write your work down.} Scratch paper is provided at the test;
  use it here too.
\item \textbf{Time yourself.} Each practice set shows a target time. Don't
  worry about speed at first---accuracy comes first, speed follows.
\item \textbf{Redo your misses a week later} and log them in the Progress
  Tracker at the back of the book.
\end{itemize}

%% ---------------------------------------------------------------- the test
\chapter*{The ASVAB Math Subtests}
\addcontentsline{toc}{chapter}{The ASVAB Math Subtests}
\markboth{The ASVAB Math Subtests}{The ASVAB Math Subtests}

The Armed Services Vocational Aptitude Battery (ASVAB) is the entrance test
for the U.S.\ Army, Navy, Air Force, Marine Corps, Space Force, and Coast Guard.
It consists of several subtests. Two of them are math:

\begin{concept}{Arithmetic Reasoning (AR)}
Word problems. Each question describes a real-life situation---a purchase,
a trip, a paycheck, a recipe, a work schedule---and you have to decide which
operations to use. Typical topics: percents and discounts, ratios, rates and
distance, averages, interest, work, measurement, and multi-step reasoning.
\end{concept}

\begin{concept}{Mathematics Knowledge (MK)}
Direct math questions: number properties, fractions, exponents and roots,
algebra (expressions, equations, inequalities, factoring), and geometry
(angles, triangles, area, circles, volume, coordinates). Knowing the
vocabulary and formulas is half the battle.
\end{concept}

\section*{Format and timing}
The ASVAB is given in two main formats. On the \textbf{computer-adaptive}
version (CAT-ASVAB), the difficulty of each question depends on how you
answered the previous ones, and you cannot go back to change an answer. On the
\textbf{paper-and-pencil} version, everyone answers the same questions and
you can move around within a subtest. Typical question counts and time
limits:\footnote{These figures can change; confirm them with your recruiter
or at officialasvab.com.}

\begin{center}
\renewcommand{\arraystretch}{1.25}
\begin{tabular}{lcc}
\toprule
 & \textbf{Arithmetic Reasoning} & \textbf{Mathematics Knowledge}\\
\midrule
CAT-ASVAB (computer) & 15 questions, 55 minutes & 15 questions, 31 minutes\\
Paper and pencil & 30 questions, 36 minutes & 25 questions, 24 minutes\\
\bottomrule
\end{tabular}
\end{center}

Every question is multiple choice with four answer choices. \textbf{Calculators
are not allowed}; scratch paper is provided. Formulas are \emph{not} given, so
memorize the ones on the Formula Sheet at the back of this book.

\section*{Why math matters so much}
Your Armed Forces Qualification Test (AFQT) score---the score that decides
whether you can enlist at all---is computed from four subtests: Arithmetic
Reasoning, Mathematics Knowledge, Word Knowledge, and Paragraph Comprehension.
The two math subtests make up roughly \textbf{half} of the AFQT. The AFQT is
reported as a percentile from 1 to 99, and each service sets its own minimum.

Arithmetic Reasoning and Mathematics Knowledge also feed many of the
\emph{line scores} (composite scores) that qualify you for specific jobs, such
as technical, electronics, and administrative specialties. A few extra correct
math answers can open up a wider choice of careers and enlistment bonuses.

%% ---------------------------------------------------------------- strategy
\chapter*{Test-Day Strategies}
\addcontentsline{toc}{chapter}{Test-Day Strategies}
\markboth{Test-Day Strategies}{Test-Day Strategies}

\begin{concept}{1. Read the last sentence first}
In a word problem, the last sentence tells you what to find. Read it first,
then read the whole problem looking for the numbers you need. Many wrong
answer choices are the right answer to a \emph{different} question (the
discount instead of the sale price, the number who passed instead of the
number who failed).
\end{concept}

\begin{concept}{2. Estimate before you calculate}
A rough answer tells you which choices are impossible. If a \$48 item is
25\% off, the sale price must be a bit more than \$35---so \$12 and \$60 are
out before you do any exact arithmetic.
\end{concept}

\begin{concept}{3. Work backward from the choices}
When a question asks ``which number\ldots'' or gives an equation, you can test
the answer choices. Start with a middle value; if it is too big, the answer is
smaller, and vice versa. For ``which value satisfies'' questions this is often
the fastest method.
\end{concept}

\begin{concept}{4. Pick numbers}
If the problem uses variables or percents of an unknown amount, choose an easy
number (100 for percents, a common multiple for fractions) and see which
answer choice gives the same result.
\end{concept}

\begin{concept}{5. Watch the units}
Minutes vs.\ hours, inches vs.\ feet, cents vs.\ dollars, radius vs.\
diameter. Convert everything to the units the question asks for
\emph{before} you calculate.
\end{concept}

\begin{concept}{6. Manage your time}
On the paper test you have a little over a minute per question. If a question
takes more than two minutes, make your best guess, mark it, and move on. Never
leave a question blank: a blank is always wrong, a guess might be right, and
on the computer test unanswered questions at the end of a subtest lower your
score.
\end{concept}

\begin{concept}{7. Check that your answer makes sense}
A sale price cannot be higher than the original price. A probability cannot
be more than 1. An average must lie between the smallest and largest values.
A quick reasonableness check catches many careless errors.
\end{concept}

\begin{tip}
Learn by heart: the squares from $1^2$ to $15^2$, the fraction--decimal--percent
equivalents ($\frac18 = 0.125 = 12.5\%$, $\frac13 \approx 33.3\%$, \ldots), and
the formulas on the Formula Sheet. Instant recall of these saves minutes on
test day.
\end{tip}

\section*{Mental Math Toolkit}
Because the ASVAB bans calculators, a handful of shortcuts pay off on almost
every page of this book. Practice them until they are automatic.

\begin{center}
\renewcommand{\arraystretch}{1.35}
\begin{tabularx}{\linewidth}{@{}>{\bfseries}p{1.55in}X@{}}
\toprule
Trick & Example\\
\midrule
Multiply by 5 & Multiply by 10, then halve: $48 \times 5 = 480 \div 2 = 240$.\\
Multiply by 25 & Multiply by 100, then divide by 4: $36 \times 25 = 3{,}600 \div 4 = 900$.\\
Multiply by 9 or 11 & $47 \times 9 = 470 - 47 = 423$; \ $34 \times 11 = 340 + 34 = 374$.\\
Break a number apart & $7 \times 98 = 7 \times 100 - 7 \times 2 = 700 - 14 = 686$.\\
Square a number ending in 5 & $35^2$: $3 \times 4 = 12$, then write 25: $1{,}225$.\\
Divide by 5 & Double, then divide by 10: $340 \div 5 = 680 \div 10 = 68$.\\
Divide by a decimal & Move both decimal points: $4.8 \div 0.06 = 480 \div 6 = 80$.\\
Build percents from 10\% & $15\%$ of $60$: $10\% = 6$, $5\% = 3$, total $9$.\\
Swap the percent & $8\%$ of $50$ = $50\%$ of $8$ = $4$.\\
Compare fractions & Cross-multiply: $\frac58$ vs.\ $\frac35$: $5 \times 5 = 25 > 8 \times 3 = 24$, so $\frac58$ is larger.\\
Estimate first & $398 \times 51 \approx 400 \times 50 = 20{,}000$ rules out answers near $2{,}000$ or $200{,}000$.\\
\bottomrule
\end{tabularx}
\end{center}
